% Generated draft: requires mathematical and source review.
\documentclass[11pt]{article}
\usepackage[margin=25mm]{geometry}
\usepackage{fontspec}
\setmainfont{Latin Modern Roman}
\usepackage{amsmath,amssymb,mathtools,unicode-math}
\setmathfont{Latin Modern Math}
\usepackage{enumitem,needspace,xurl,hyperref}
\setlength{\parindent}{0pt}
\setlength{\parskip}{5pt}
\setlength{\emergencystretch}{3em}
\newcommand{\N}{\mathbb{N}}
\newcommand{\Z}{\mathbb{Z}}
\newcommand{\Q}{\mathbb{Q}}
\newcommand{\R}{\mathbb{R}}
\newcommand{\C}{\mathbb{C}}
\newcommand{\F}{\mathbb{F}}
\newcommand{\A}{\mathbb{A}}
\newcommand{\PP}{\mathbb{P}}
\newcommand{\E}{\mathbb{E}}
\newcommand{\eps}{\varepsilon}
\newcommand{\dd}{\,\mathrm{d}}
\begin{document}
\section*{UnsolvedMath ID: 800398}
\textbf{Source problem code:} catalogue entry
\subsection{Definitions and mathematical statement}
Let \(d\geq 2\) be an integer and let \(\mathcal P\) be a portrait in the sense intended by the source. For a number field \(K\) and a finite set \(S\) of primes of \(K\), write \(R_S\) for the ring of \(S\)-integers, with the usual convention
\[
R_S=\{x\in K: v(x)\geq 0\text{ for every nonarchimedean prime }v\notin S\}.
\]
Retain the source's notation
\[
Q_{d,\mathcal P}(K,S):=
\mathcal{G\!R}_d^1[\mathcal P](K,S)/\operatorname{PGL}_2(R_S).
\]
Here the supplied excerpt does not define a portrait, the collection \(\mathcal{G\!R}_d^1[\mathcal P](K,S)\), or the action used to form this quotient. Consequently, this notation records the intended source object without supplying its missing definition; in particular, no particular good-reduction or portrait-realization conditions are being assumed.

The requested classification is to give criteria depending on \(d\) and the geometry of \(\mathcal P\) that determine which of the following holds:
\begin{enumerate}
\item For every number field \(K\) and every finite set \(S\) of primes of \(K\),
\[
\#Q_{d,\mathcal P}(K,S)<\infty.
\]
\item There exist a number field \(K\) and a finite set \(S\) of primes of \(K\) such that
\[
\#Q_{d,\mathcal P}(K,S)=\infty.
\]
\end{enumerate}
These alternatives are complementary once the source objects and their domains have been specified. A fully defined mathematical formulation requires the source's definitions of the allowed portraits, their geometry, the collection above, and its group action. The supplied excerpt also does not specify whether sets \(S\) must contain the archimedean primes; those primes impose no conditions in the displayed definition of \(R_S\), but any additional convention in the source remains unresolved.
\subsection{Short English statement}
Determine, from the degree \(d\) and the geometry of a portrait \(\mathcal P\), whether the source's quotient is finite over every number field and every allowed finite set of primes, or infinite for at least one such pair. The excerpt omits definitions needed to identify the objects being classified.
\subsection{Sources}
\begin{itemize}
\item Local UnsolvedMath registry, permanent ID \(800398\), exact source problem code \texttt{catalogue entry}, title ``Open Problems in Arithmetic Dynamics --- Question 21.'' The supplied record attributes the question to Benedetto, DeMarco, Ingram, Jones et al., with proposed year 2018.
\item Benedetto, DeMarco, Ingram, Jones et al., \emph{Current Trends and Open Problems in Arithmetic Dynamics} (2018), Question 21. Supplied source URL: \url{https://arxiv.org/abs/1806.04980}. The background records access on 2026-10-16 and extraction method \texttt{arxiv\_latex}; these are supplied provenance claims, not an independent source visit or verification.
\item The local record labels the problem ``open'' and records its update at \texttt{2026-10-16T14:00:00Z}. This is supplied registry curation, not an independently verified assessment of the mathematical literature or present status.
\end{itemize}
\end{document}
